% Propietario reproducido por unidad demostrativa; véase PROPIETARIOS_CONTINUO_UNIVERSO.json.
% Origen: XI ES CUMPLIMIENTO_EDITORIAL_20260928, sections/ch_cierre_cardinal.tex:26--403.
% Cuerpos y pruebas conservados; sólo namespace, rutas y remisiones contextuales declaradas.
\subsection{Prolongación de estados y extensión semántica transfinita}

Sea

\[
 \Sigma_{\rm HMT}
 =\bigl(\dot R_{\rm APP},\dot R_{\rm TRIT},\dot R_{\rm TPK},
 \dot R_{U},\dot R_{\Gamma_9},\dot R_{\rm Ext},\dot R_\rho,
 \dot R_{\pi},\dot R_{\varphi},\dot R_e,\dot R_{\alpha},
 \dot R_{\rm ar},\dot R_{\rm exc}\bigr)
\tag{33}\label{hmtfg:base:eq:ch-signatura}
\]

la signatura \textbf{relacional} numerable cuyo código efectivo es \(h\). Cada
\(\dot R_S\) es el grafo del operador \(S\) sobre los códigos donde está
definido; no se presupone que una función externa sea total dentro de cada
nivel. Los grafos \(\dot R_{\pi},\dot R_{\varphi},\dot R_e,\dot R_{\alpha}\)
son los lectores de salida construidos por la generación coinductiva y el
cierre dodecafásico; no reciben como argumentos los valores convencionales de
las constantes. Su presencia en \(\Sigma_{\rm HMT}\) asegura que la clausura
conserve el estado completo: integra simultáneamente las rutas que expresan
valor y la incidencia que se prolonga hacia la realización excepcional.

En cada horizonte, el código del estado registra las coordenadas APP--TRIT--TPK,
la ruta, la hoja, el acarreo, la frontera y la memoria; en sus etapas de
representación registra asimismo \(P_N,\Phi_N,E_N,A_N\), donde
\(A_N\) es el prefijo de \(\alpha_{\rm HMT}\) producido por el mismo estado.
Estas coordenadas no seleccionan la rama desde fuera: son parte de la historia
que se codifica. Fijada una numeración primitivo-recursiva de las coordenadas finitas de
un estado, toda historia inversa
\(x_\infty=(x_N)_{N\geq1}\in X_\infty^{\rm enr}\) determina el código
\[
 \operatorname{code}(x_\infty)
 :=\{\langle N,j,c\rangle:(x_N)_j\text{ tiene código }c\}\subseteq\omega.
\]
Las ecuaciones de coherencia
\(\rho_{N+1,N}(x_{N+1})=x_N\) permiten decodificar de manera única la
historia a partir de ese real. Escribimos \(m_\infty\) para el código de la
historia realizada y \(m_\infty(n)\) para su función característica bajo la
numeración fijada. El código operatorio \(h\) y este registro completo se
reúnen, mediante una función primitiva recursiva de emparejamiento, en el real

\[
 u_{h,m}:=
 \{\langle0,\langle n,h(n)\rangle\rangle:n<\omega\}
 \cup
 \{\langle1,\langle n,m_\infty(n)\rangle\rangle:n<\omega\}.
\tag{33a}\label{hmtfg:base:eq:ch-codigo-conjunto}
\]

Las dos proyecciones primitivas recuperan \(h\) y \(m_\infty\) desde
\(u_{h,m}\).

La primera extensión es la prolongación métrica y operatoria de los estados
finitos. Para \(N\ge1\), sea
\(\widehat{\mathcal X}^{\rm enr,raw}_{6N}\) el conjunto de estados
enriquecidos brutos de longitud \(6N\), y defínase conjuntistamente el
subsistema superviviente por
\[
\begin{aligned}
 X_N
 &:=\widehat{\mathcal X}^{\rm enr,surv}_{6N}\\
 &:=\Bigl\{x\in\widehat{\mathcal X}^{\rm enr,raw}_{6N}:\\[-.2ex]
 &\qquad\forall M>N\ \exists z\in
   \widehat{\mathcal X}^{\rm enr,raw}_{6M}\;\text{ tal que}\\[-.2ex]
 &\qquad\operatorname{tr}_{6M,6N}(z)=x\ \wedge\
   z\text{ satisface toda transición TPK intermedia}\Bigr\}.
\end{aligned}
\]
y póngase
\[
 \rho_{N+1,N}:=\operatorname{tr}_{6N+6,6N}\!\upharpoonright X_{N+1},
\]
y
\[
 \operatorname{Ext}_N(x):=
 \{y\in X_{N+1}:(x,y)\in
 \operatorname{Ext}_{6N+6,6N}\}.
\]
El nivel raíz se incluye sin excepción implícita:
\[
 X_0:=\{\ast\},\qquad
 \rho_{1,0}:X_1\to X_0\text{ es constante},\qquad
 \operatorname{Ext}_0(\ast):=X_1.
\]
El superíndice \(\mathrm{surv}\) denota el
árbol podado de los estados que poseen descendientes TPK compatibles a todo
horizonte finito. La ley de prolongación TPK establecida en las secciones
precedentes prueba que las semillas
emitidas pertenecen a ese árbol; la ramificación finita y el lema de Kőnig
convierten la compatibilidad a todo horizonte en una rama infinita. En
particular,

\[
 \forall N\ge0\;\forall x\in X_N\quad
 \varnothing\ne\operatorname{Ext}_N(x)
 \subseteq\rho_{N+1,N}^{-1}(x).
\tag{34}\label{hmtfg:base:eq:ch-extension-no-vacia}
\]

Por el mismo lema de árbol finitamente ramificado, toda semilla superviviente
posee una prolongación compatible

\[
 x_\infty=(x_N)_{N\ge1}\in
 X_\infty^{\rm enr}:=\varprojlim_N(X_N,\rho_{N+1,N}).
\tag{35}\label{hmtfg:base:eq:ch-limite-inverso}
\]

La segunda extensión actúa \textbf{después} sobre el código completo de esa
historia realizada.
No se infiere de la mera existencia del operador finito
\(\operatorname{Ext}_N\): formaliza a escala conjuntista el principio HMT de
que sólo son objetos del dominio íntegro los que poseen una derivación
genealógica desde el estado y sus reglas. La extensión semántica transfinita
es:

\[
 \begin{aligned}
  \mathfrak G_0(h,m_\infty)
  &:=\operatorname{TC}\bigl(\{\omega,u_{h,m},h,m_\infty\}\bigr),\\
  \mathfrak G_{\beta+1}(h,m_\infty)
  &:=\operatorname{Def}_{\Sigma_{\rm HMT}}
       \bigl(\mathfrak G_\beta(h,m_\infty)\bigr),\\
  \mathfrak G_\lambda(h,m_\infty)
  &:=\bigcup_{\beta<\lambda}\mathfrak G_\beta(h,m_\infty)
       \quad(\lambda\text{ límite}),\\
  \mathfrak G_{\rm HMT}(h,m_\infty)
  &:=\bigcup_{\beta\in\operatorname{Ord}}
       \mathfrak G_\beta(h,m_\infty).
 \end{aligned}
\tag{36}\label{hmtfg:base:eq:ch-clausura-transfinita}
\]

La hipótesis del continuo y una biyección objetivo no se introducen entre las
condiciones que definen el operador sucesor; el operador está dado
prospectivamente por

\[
 \operatorname{Def}_{\Sigma_{\rm HMT}}(M)
 :=M\cup
 \Bigl\{
   \{x\in M:\langle M,\in,
       (\dot R_S\cap M^{k_S})_{S\in\Sigma_{\rm HMT}}\rangle
       \models\varphi(x,\bar p)\}:
   \ulcorner\varphi\urcorner\in\omega,
   \ \bar p\in M^{<\omega}
 \Bigr\}.
\tag{37}\label{hmtfg:base:eq:ch-operador-def}
\]

Es decir: expresa simultáneamente todos los subconjuntos definibles por una
fórmula finita del lenguaje relacional HMT y parámetros que la genealogía ya
produjo. No
consulta un valor real objetivo, una constante convencional, un cardinal
objetivo ni un modelo exterior. La fórmula~\eqref{hmtfg:base:eq:ch-operador-def} define el cierre transfinito
sobre las semillas, los operadores, el estado enriquecido, el registro y la
doble proyección construidos previamente.

\begin{lemma}[Codificación conjunta del programa HMT y del registro realizado]
\label{hmtfg:base:thm:ch-6}
El par formado por
el código operatorio \(h\) y el registro \(m_\infty\) se codifica por un único
real \(u_{h,m}\subseteq\omega\), y toda fórmula de \(\Sigma_{\rm HMT}\) se
traduce uniformemente a una fórmula del lenguaje \(\{\in,u_{h,m}\}\).

\end{lemma}
\begin{proof}
Tanto \(h\) como \(m_\infty\) son sucesiones de naturales: \(h\)
codifica la signatura, las tablas finitas y los programas de los operadores;
\(m_\infty\) codifica una sola historia realizada, no el espacio completo de
historias. La fórmula~\eqref{hmtfg:base:eq:ch-codigo-conjunto} y sus dos decodificadores prueban la primera
afirmación. Cada símbolo de~\eqref{hmtfg:base:eq:ch-signatura} designa su \textbf{relación de grafo codificada en \(h\)}, nunca el grafo externo de todas sus posibles evaluaciones. APP,
TRIT, \(U_t\), \(\Gamma_9\), las restricciones \(\rho\) y las extensiones
finitas tienen grafos primitivo-recursivos sobre códigos enteros y, por tanto,
fórmulas \(\Delta_0\) absolutas entre niveles transitivos. Las dos
representaciones se expresan como relaciones de grafo mediante las fórmulas

\[
 \begin{aligned}
 \dot R_{\rm ar}(x,y)
 &\;\Longleftrightarrow\;
 \forall k\in\omega\;\exists N\;\forall n\ge N\quad
 |\mu_n(x)-y|<2^{-k},\\
 \dot R_{\rm exc}(x,c,z)
 &\;\Longleftrightarrow\;
 \forall N\in\omega\quad z\!\upharpoonright N
   =E_N(x\!\upharpoonright N,c\!\upharpoonright N),
 \end{aligned}
\tag{37c$_0$}\label{hmtfg:base:eq:ch-grafos-publicacion}
\]

donde \(\mu_n\) y \(E_N\) son las funciones racionales/finitas cuyos códigos
figuran en \(h\). Se sustituye inductivamente cada fórmula atómica que contiene
un símbolo HMT por su fórmula de grafo con parámetro \(u_{h,m}\). Los
conectivos y cuantificadores se conservan. Esta inducción sobre la complejidad
sintáctica produce una traducción \(\varphi\mapsto\varphi^*\) tal que, para
todo nivel transitivo \(M\) de~\eqref{hmtfg:base:eq:ch-clausura-transfinita} que contiene los parámetros,

\[
 \langle M,\in,(\dot R_S\cap M^{k_S})_S\rangle
   \models\varphi(\bar a)
 \quad\Longleftrightarrow\quad
 M\models\varphi^*(\bar a;u_{h,m}).
\tag{37c}\label{hmtfg:base:eq:ch-eliminacion-uniforme}
\]

Así, \eqref{hmtfg:base:eq:ch-operador-def} no puede consultar un oráculo, una tabla decimal objetivo ni una
relación no codificada que introduzca incontables objetos en un solo paso.
\end{proof}

La semántica HMT se define ahora \textbf{independientemente de identificarla con \(\mathfrak G\)}. Sea \(\mathcal T_0^{\rm HMT}\) el conjunto de nombres
canónicos de los elementos de
\(\mathfrak G_0(h,m_\infty)=
\operatorname{TC}(\{\omega,u_{h,m},h,m_\infty\})\). En particular, el nivel
inicial contiene el registro realizado \(m_\infty\), no todas las historias de
\(X_\infty^{\rm enr}\) como parámetros primitivos.
Los términos admisibles son los árboles bien fundados obtenidos por

\[
 \begin{aligned}
 \mathcal T^{\rm HMT}_{\beta+1}
 &:={\mathcal T}^{\rm HMT}_\beta\cup
 \{\langle\beta,\ulcorner\varphi\urcorner,
       t_0,\ldots,t_{k-1}\rangle:
       t_i\in\mathcal T^{\rm HMT}_\beta\},\\
 \mathcal T^{\rm HMT}_\lambda&:=
       \bigcup_{\beta<\lambda}\mathcal T^{\rm HMT}_\beta,
 \qquad
 \mathcal T^{\rm HMT}:=
       \bigcup_{\beta\in\operatorname{Ord}}\mathcal T^{\rm HMT}_\beta.
 \end{aligned}
\tag{37a}\label{hmtfg:base:eq:ch-terminos}
\]

El valor de un término sucesor es el subconjunto del nivel anterior definido
por \(\varphi\) en la estructura relacional inducida, con los valores de
\(t_0,\ldots,t_{k-1}\) como parámetros. Se define, antes de cualquier
comparación conjuntista,

\[
 \operatorname{Ob}_{\rm sem}^{\rm HMT}(h,m_\infty)
 :=\{\llbracket t\rrbracket:t\in\mathcal T^{\rm HMT}\}.
\tag{37b}\label{hmtfg:base:eq:ch-objetos-semanticos}
\]

Ésta es la regla de \textbf{exhaustividad generativa} adoptada por HMT: todo
objeto del dominio íntegro debe poseer uno de esos árboles de derivación y todo
árbol admisible expresa un objeto. No es consecuencia de la ramificación
finita. La hipótesis del continuo no se introduce en esta regla; se obtiene
mediante la demostración desarrollada en el artículo. Tampoco se introducen
como datos una enumeración de los reales ni una biyección objetivo.

\begin{theorem}[Equivalencia entre la semántica de términos y la jerarquía genealógica]
\label{hmtfg:base:thm:ch-6a}
La evaluación de
los términos anteriores satisface

\[
 \operatorname{Ob}_{\rm sem}^{\rm HMT}(h,m_\infty)
 =\mathfrak G_{\rm HMT}(h,m_\infty).
\tag{37d}\label{hmtfg:base:eq:ch-equivalencia-semantica}
\]

\end{theorem}
\begin{proof}
Por inducción transfinita. En el nivel inicial las dos clases
coinciden. Si todo valor de un término de
\(\mathcal T^{\rm HMT}_\beta\) pertenece a
\(\mathfrak G_\beta\), la semántica del término sucesor es uno de los
subconjuntos de~\eqref{hmtfg:base:eq:ch-operador-def}. Recíprocamente, cada subconjunto de~\eqref{hmtfg:base:eq:ch-operador-def} viene dado por
un código de fórmula y una tupla finita de parámetros; por la hipótesis
inductiva esos parámetros poseen términos y~\eqref{hmtfg:base:eq:ch-terminos} forma el término que lo
denota. En un límite ambas construcciones toman la misma unión.
\end{proof}

\Needspace{9\baselineskip}
El codificador de estados es compatible con los truncamientos y envía cada
prolongación TPK a un término sucesor. Sea
\(\mathcal T_N^{\rm st}\subset\mathcal T^{\rm HMT}\) la subclase de términos
que codifica estados completos de profundidad \(N\), y sea
\(\tau_{N+1,N}:\mathcal T_{N+1}^{\rm st}\to\mathcal T_N^{\rm st}\) el borrado
del último bloque junto con el truncamiento correspondiente del registro genealógico,
acarreo y memoria. Si \(J_N:X_N\to\mathcal T_N^{\rm st}\) es el codificador
completo, la compatibilidad de \(\operatorname{Ext}_N\) prueba la identidad
correctamente tipada

\[
 \forall x\in X_N\ \forall y\in\operatorname{Ext}_N(x)\qquad
 \tau_{N+1,N}\bigl(J_{N+1}(y)\bigr)=J_N(x).
\tag{37e$_0$}\label{hmtfg:base:eq:ch-naturalidad-codificador}
\]

Defínase dentro de la semántica
\[
 \operatorname{Succ}^{\rm TPK}_{N+1}(J_N(x)):=
 \{t\in\mathcal T_{N+1}^{\rm st}:\exists y\in X_{N+1}\,
   (y\in\operatorname{Ext}_N(x)\ \wedge\ t=J_{N+1}(y))\}.
\]
Elíjase \(\beta\) tal que \(\mathfrak G_\beta\) contenga \(x\), \(X_N\),
\(X_{N+1}\), los grafos finitos de
\(\operatorname{Ext}_N,J_N,J_{N+1}\) y sus códigos en \(h\). Entonces la
fórmula primitivo-recursiva del grafo de \(\operatorname{Ext}_N\) da

\[
 J_{N+1}[\operatorname{Ext}_N(x)]
 =\operatorname{Succ}^{\rm TPK}_{N+1}(J_N(x)),
 \qquad
 J_{N+1}[\operatorname{Ext}_N(x)]
 \in\mathfrak G_{\beta+1},
\tag{37e}\label{hmtfg:base:eq:ch-sucesores-semanticos}
\]

y ambos miembros son formados por la fórmula de grafo de la misma extensión
TPK. Las ecuaciones~\eqref{hmtfg:base:eq:ch-naturalidad-codificador}--\eqref{hmtfg:base:eq:ch-sucesores-semanticos} son el enlace semántico explícito: la
subclase de términos sucesores es exactamente la imagen de las prolongaciones
admisibles y su truncamiento conmuta. No se identifica la fibra completa de
\(\rho\) con \(\operatorname{Ext}_N\), ni se identifica
\(\operatorname{Def}\) con TPK por nomenclatura.

\begin{lemma}[Transitividad, axiomas y buen orden del modelo genealógico]
\label{hmtfg:base:thm:ch-6b}
\(\mathfrak G_{\rm HMT}(h,m_\infty)\) es una clase transitiva, contiene todos
los ordinales, satisface la teoría de conjuntos de Zermelo--Fraenkel (ZF) y
posee un buen orden global definible; satisface,
por tanto, elección y todos los axiomas usados en la comparación cardinal posterior.

\end{lemma}
\begin{proof}
Abreviemos \(B=\mathfrak G_0\) y \(u=u_{h,m}\). La inducción
sobre~\eqref{hmtfg:base:eq:ch-clausura-transfinita} prueba simultáneamente que cada nivel es transitivo y que
\(\mathfrak G_\alpha\in\mathfrak G_{\alpha+1}\); si todos los ordinales
menores que \(\alpha\) están presentes, su conjunto es definible en el nivel
correspondiente y expresa \(\alpha\). Por ello la unión contiene todos los
ordinales. Extensionalidad y Fundación se heredan del universo ambiente;
Vacío e Infinito están en \(B\). Elegidos un nivel que contenga los parámetros,
las fórmulas que definen \(\{a,b\}\), \(\bigcup a\) y los subconjuntos por
Separación los expresan en un sucesor.

Para Colección--Reemplazo, sea \(F\) una relación funcional definible en
\(\mathfrak G\) sobre el conjunto \(a\). El esquema de reflexión relativo a
la jerarquía definible \((\mathfrak G_\xi,\in,u_{h,m})_{\xi\in\mathrm{Ord}}\),
aplicado en la metateoría ambiente a la lista finita de fórmulas que define
\(F\), produce un nivel
\(\mathfrak G_\theta\) que contiene \(a\), los parámetros y es correcto para
esas fórmulas. Los rangos de los testigos de \(F[a]\) están entonces acotados
por \(\theta\); Separación sobre \(\mathfrak G_\theta\) forma la imagen en
\(\mathfrak G_{\theta+1}\). Para Potencia no se presupone el conjunto interno
que se quiere construir. Fijado \(a\in\mathfrak G_\gamma\), se considera en
la metateoría ambiente el conjunto ya existente \(\mathcal P^V(a)\). Para
cada \(b\in\mathcal P^V(a)\cap\mathfrak G\), sea \(r(b)\) su primera etapa
genealógica. Como \(\mathfrak G\) es una clase definible, Separación en la
metateoría ambiente forma primero el conjunto
\(\mathcal P^V(a)\cap\mathfrak G\). Reemplazo \textbf{en la metateoría ambiente}
aplicado a ese conjunto acota los ordinales \(r(b)\) por algún \(\theta\).
En consecuencia,
\[
 \mathcal P^{\mathfrak G}(a)
 =\{b\in\mathfrak G_\theta:b\subseteq a\},
\]
y el miembro derecho se expresa por Separación en
\(\mathfrak G_{\theta+1}\). No introduce en \(\mathfrak G\) los subconjuntos
exteriores a la semántica~\eqref{hmtfg:base:eq:ch-objetos-semanticos}. Éste es el argumento no circular de
acotación por rangos.

Finalmente, a cada objeto se le asigna su primer nivel de nacimiento y su
menor \textbf{nombre de derivación}: un árbol de derivación bien fundado y
finitamente ramificado (y, por~\eqref{hmtfg:base:eq:ch-terminos}, finito para cada nombre individual), etiquetado por un
ordinal de nivel, un código de fórmula y la tupla de nombres de sus parámetros.
Por recursión transfinita se ordenan esos nombres: primero la etiqueta ordinal,
después el código de fórmula y después, lexicográficamente, los nombres de los
parámetros ya ordenados. El menor nombre de cada valor define una relación de
clase \(<_{\mathfrak G}\) que compara todos los objetos. Es un buen orden
global definible y, tomando el menor elemento de cada conjunto no vacío,
prueba Elección.
\end{proof}

% Propietario reproducido por unidad demostrativa; véase PROPIETARIOS_CONTINUO_UNIVERSO.json.
% Origen: XI ES CUMPLIMIENTO_EDITORIAL_20260928, sections/ch_reflexion_testigos.tex:1--89.
% Cuerpos y pruebas conservados; sólo namespace, rutas y remisiones contextuales declaradas.
\paragraph{Reflexión por clausura de testigos y construcción de la imagen.}
El paso de Colección--Reemplazo del lema~\ref{hmtfg:base:thm:ch-6b} admite la siguiente
demostración explícita. Se mantiene la jerarquía ya generada a partir del
código operatorio y del registro realizado; abreviamos
\(\mathfrak G=\mathfrak G_{\rm HMT}(h,m_\infty)\).
La argumentación utiliza Separación y Reemplazo en la metateoría ambiente,
sin presuponer estos esquemas en \(\mathfrak G\).

Sea \(\Phi\) una colección finita de fórmulas. Se traducen primero sus
cuantificadores universales mediante negación y cuantificación existencial,
y se cierra la colección resultante bajo subfórmulas. Para cualquier nivel
inicial existe un ordinal límite posterior \(\theta\) tal que
\[
 \mathfrak G_\theta\models\varphi(\bar a)
 \quad\Longleftrightarrow\quad
 \mathfrak G\models\varphi(\bar a)
 \qquad
 (\varphi\in\Phi,\ \bar a\in\mathfrak G_\theta^{<\omega}).
\]
La notación de satisfacción de la clase se interpreta, para cada fórmula
fijada, relativizando sus cuantificadores a la clase definible
\(\mathfrak G\); no se introduce un predicado uniforme de verdad.

\begin{proof}
Fijemos \(\alpha\). Para cada subfórmula existencial
\(\exists y\,\psi(y,\bar x)\) de \(\Phi\), las tuplas
\(\bar a\in\mathfrak G_\alpha^{<\omega}\) de la aridad correspondiente
que poseen un testigo en \(\mathfrak G\) forman un conjunto por Separación
metateórica. A cada una se asigna el menor ordinal \(\beta\) para el que
algún \(y\in\mathfrak G_\beta\) satisface
\(\mathfrak G\models\psi(y,\bar a)\). El ordinal existe porque
\(\mathfrak G\) es la unión de sus niveles. Reemplazo en la metateoría
acota esas primeras etapas de testigos. La finitud de \(\Phi\) permite
elegir definiblemente una sola cota \(f(\alpha)>\alpha\), válida para
todas sus subfórmulas existenciales. Por tanto, cada tupla considerada tiene
\emph{al menos un testigo} en \(\mathfrak G_{f(\alpha)}\).

Tomamos \(\alpha_0\) suficientemente grande para contener los parámetros
iniciales y ponemos
\[
 \alpha_{n+1}=f(\alpha_n),\qquad
 \theta=\sup_{n<\omega}\alpha_n.
\]
Por continuidad de la jerarquía, toda tupla finita de
\(\mathfrak G_\theta\) pertenece a algún \(\mathfrak G_{\alpha_n}\).
Si una fórmula existencial de \(\Phi\) con esa tupla es verdadera en
\(\mathfrak G\), posee un testigo en
\(\mathfrak G_{\alpha_{n+1}}\subseteq\mathfrak G_\theta\).
La inducción sobre las fórmulas prueba ahora la equivalencia: las atómicas
se conservan en la estructura inducida; los conectivos preservan la
equivalencia; el paso existencial utiliza ese testigo en una dirección y el
mismo testigo, visto en \(\mathfrak G\), en la otra. La traducción previa
de los universales incluye también las fórmulas existenciales que expresan
sus posibles contraejemplos. Queda probada la reflexión finita requerida.
\end{proof}

Sea ahora \(F(x,y,\bar p)\) funcional sobre \(a\) en \(\mathfrak G\).
Aplicamos la construcción a las fórmulas que expresan \(F\), la existencia
de sus testigos y su imagen, con \(a,\bar p\in\mathfrak G_\theta\).
La transitividad da \(a\subseteq\mathfrak G_\theta\). Cada
\(x\in a\) posee su valor dentro de ese nivel, y la reflexión identifica
la imagen exacta con
\[
 b=\{y\in\mathfrak G_\theta:
       \mathfrak G_\theta\models
       \exists x\in a\ F(x,y,\bar p)\}.
\]
Este conjunto es definible sobre \(\mathfrak G_\theta\), con parámetros
del mismo nivel. Por la operación sucesora de la jerarquía,
\(b\in\operatorname{Def}(\mathfrak G_\theta)
\subseteq\mathfrak G_{\theta+1}\). Así se obtiene Reemplazo en
\(\mathfrak G\). Si se conserva la existencia de testigos y se omite la
unicidad, la misma construcción proporciona un conjunto de testigos para
Colección.

\paragraph{Aridad de los nombres y selección interna.}
En el buen orden del lema~\ref{hmtfg:base:thm:ch-6b}, la enumeración de la base
asigna a cada elemento su menor código natural. En cada sucesor se ordenan
primero los códigos de fórmula; fijado uno, queda fijada la aridad de su
tupla de parámetros. El orden lexicográfico se aplica entonces a un
\emph{producto finito de aridad fija} de los buenos órdenes anteriores.
El menor nombre representa cada valor nuevo y los niveles límite conservan
el orden por primera aparición. De este modo queda especificado el orden
de las tuplas, sin mezclar aridades distintas en un único orden
lexicográfico. La selección del menor elemento de cada conjunto no vacío,
junto con Reemplazo ya obtenido, forma internamente las funciones de
elección utilizadas después. Esta construcción de nombres precede a las
inyecciones cardinales y no recibe ninguna biyección del continuo como
dato.

